\chapter{Experimental Setup}

\section{Data}

All training runs use prompts from DeepMath-103K \cite{he2025DeepMath103KLargeScaleChallenging}, a large-scale mathematical reasoning dataset available on Hugging Face\footnote{https://huggingface.co/datasets/zwhe99/DeepMath-103K}. This dataset is well matched to our target setting because it contains difficult, verifiable mathematics problems rather than general instruction-following prompts. The problems span a broad range of mathematical topics and were designed to be challenging and decontaminated from common evaluation benchmarks. We use the prompts as the common input distribution across all methods so that differences between runs can be attributed to the distillation objective and teacher choice rather than to differences in task selection.

For off-policy experiments that require precomputed teacher data rollouts, we generate the teacher completions with DataTrove \cite{penedo2024datatrove} using the hyperparameters described in Appendix~\ref{app: hyperparameters}. We make the generated datasets available on Hugging Face.

\section{Models}

We use the Qwen3 model family \cite{yangQwen3TechnicalReport2025} to study different levels of the capacity gap. Keeping teacher and student within the same model family is important because architecture, tokenizer, training data, and chat-template differences can introduce stylistic and distributional mismatches that are unrelated to model size. Those mismatches can degrade distillation performance and obscure the specific effect of the teacher--student capacity gap \cite{liRethinkingOnPolicyDistillation2026}. Crucially, \cite{yangQwen3TechnicalReport2025} report that the small models go through a distillation process from larger models in the Qwen3 family. This prior distillation makes it even more important to be consistent with the model family from the teacher previously used for distillation \cite{wuLightningOPDEfficient2026}. We include preliminary experiments with different model families in Appendix~\ref{app: different_model_family}, but the main experiments focus on Qwen3.

Specifically, we use Qwen3-0.6B and Qwen3-4B in non-thinking mode as students, and Qwen3-4B-Instruct, Qwen3-30B-A3B-Instruct-2507, and Qwen3-235B-A22B-Instruct-2507 as teachers. For the initial student models, we modify only the chat template, not the weights, to force non-thinking mode. This ensures that the student remains in the intended inference mode throughout training and preserves the baseline behavior reported by the Qwen team.

\section{Inference Engine and Computational Resources} \label{sec: inference_engine}

We use vLLM \cite{kwon2023efficient} to generate student rollouts and teacher log probabilities. vLLM varies significantly between versions, so we pin version 0.17.1 for all our experiments. Using 0.17.1 instead of the latest 0.20.0 means we can only get the top-k log-probabilities from a vLLM server but cannot get specific log-probabilities. As shown in Figure~\ref{fig:topk_forward_vs_reverse_kl}, we require specific log-probabilities from the teacher based on the top-k under the student to calculate the strict top-k RKL. That is why we need to define the alternative TS-RKL formulation explained in Section \ref{sec: top-k_kl}

The 4B teacher fits on the same GPUs used for training, while the 30B and 235B teachers are served through separate vLLM servers. Throughout the experiments, student training uses 4 nodes with 8 NVIDIA H100 GPUs each. The 30B teacher uses 1 additional serving node, and the 235B teacher uses 2 additional serving nodes. We include all the setting for rollout generation and server configuration in the Appendix ~\ref{app: hyperparameters}.


\section{Benchmarks}

Our target capability is mathematical reasoning, so AIME25 \cite{aime25} is the main evaluation benchmark for the experiments. AIME is a competition-style high-school mathematics exam with short numeric answers. The problems require multi-step algebraic, geometric, combinatorial, and number-theoretic reasoning, and they are difficult to solve by pattern matching alone. This makes AIME25 a useful endpoint for measuring whether distillation improves the kind of structured mathematical reasoning targeted by DeepMath-103K.

We also evaluate on two auxiliary benchmarks to measure whether training on mathematics improves the target capability at the cost of broader behavior. GPQA \cite{reinGPQAGraduateLevelGoogleProof2023} measures graduate-level scientific reasoning through multiple-choice questions written by domain experts in biology, physics, and chemistry. The questions are designed to be difficult for skilled non-experts even with web access, so performance on GPQA gives a signal about whether the model preserves harder out-of-domain reasoning skills. IFEval \cite{zhou2023instructionfollowingevaluationlargelanguage} measures instruction following using prompts with automatically verifiable constraints, such as required keywords, formatting rules, or length restrictions. We include it because reasoning fine-tuning and distillation can change a model's compliance with explicit user instructions even when the training data is not directly about instruction following.

We run all evaluations with the Inspect library \cite{UK_AI_Security_Institute_Inspect_AI_Framework_2024}. Using the same evaluation harness across checkpoints keeps prompting, sampling, answer extraction, and scoring consistent, which is especially important when comparing small changes between distillation objectives and teacher sizes.

\section{Implementation} \label{sec: implementation}

Distillation requires access to log probabilities from both the teacher and the student in order to compute the losses above. In OPD this requirement is coupled to rollout generation: the training examples are not fixed completions, but samples from the current student policy $\pi_\theta$. For small teachers, the simplest implementation is to colocate the teacher, the student, and the rollout generator on the same training GPUs. This becomes impractical when the teacher is tens or hundreds of billions of parameters, because the training GPUs must already store the student, optimizer state, activations for backpropagation, and the inference engine used for student generation. We therefore separate the teacher from the training job and query it through an external server, following the large-teacher setup described in \cite{patino2026_distilling_100b_models_40x_faster_with_trl}.

The on-policy part of the algorithm creates a systems constraint that is absent in ordinary supervised distillation. Rollouts cannot be generated once, cached, and reused throughout training, because they must come from the current student rather than from an older checkpoint. This means generation happens inside the training loop. In our implementation, vLLM is used as the inference engine for efficient student rollouts, while Transformers \cite{wolf-etal-2020-transformers} is responsible for the student forward pass, log-probability computation, backpropagation, and optimizer updates. These two views of the student must be kept synchronized. After the optimizer changes the Transformers model weights, the vLLM copy used for generation must be refreshed before it is allowed to produce new on-policy samples. Otherwise the loss would be computed on trajectories from a stale policy, weakening the on-policy guarantee. A generation buffer can improve throughput by batching prompts across gradient accumulation steps, but only because the student weights remain fixed until the optimizer step.

Serving the teacher on separate GPUs solves the memory problem but turns teacher scoring into a distributed inference workload. In this setup, the \emph{client} is the set of GPUs and processes used for training the student, and the \emph{server} is the set of GPUs running the teacher model behind vLLM. The client sends the student-generated sequences to the server and receives the teacher log probabilities needed by the loss. This communication is expensive because log-probability responses are much larger than ordinary text-generation responses. Even with a top-$k$ approximation, the server must return token ids and log probabilities for many positions in each sequence. This is especially challenging on the reasoning tasks we are interested in, where the number of tokens is large. The second challenge is latency. Multiple data-parallel training workers can query the same teacher server concurrently, and without careful request batching the server either underutilizes vLLM's batching capacity or accumulates long queues that leave the training GPUs idle. For this reason, the teacher server must be engineered as part of the training system, with batching and compact serialization used to reduce both transfer cost and tail latency.

Our implementation addresses both of these challenges and enables distillation with large teachers. We published our implementation in the TRL library \cite{vonwerra2020trl} as the DistillationTrainer class and provide a detailed benchmarking of our methods in a separate publication \cite{patino2026_distilling_100b_models_40x_faster_with_trl}.